\section*{Formalization and Palomar registrations}
\phantomsection\label{sec:formalization}
\addcontentsline{toc}{section}{Formalization and Palomar registrations}
A companion formalization in Lean~4 covers the principal arithmetic constructions
and probabilistic conclusions of this article and its technical companion.
The development uses Lean~4.34.0 and Mathlib~v4.34.0 and is available in the
\href{https://github.com/Vulkin-prog/paper-c-lean}{\texttt{paper-c-lean} repository}.
It includes explicit definitions of the random model, conditioning events,
observables and comparison laws, together with proofs and reproducible verification
instructions. The repository's correspondence between the manuscripts and Lean
statements records the precise coverage, hypotheses and proof substitutions.

The formalization is relative to seven explicit literature inputs: scalar Poisson
Stein factors; finite Arratia--Goldstein--Gordon process comparison; the prime
number theorem remainder; multivariate Poisson Stein solution bounds; the uniform
prime-divisor estimate of Laishram--Shorey; Shorey's square-product exclusion; and
the Nicolas--Robin divisor bound. These propositions are hypotheses of the formal
theorems that use them; their proofs are not supplied by this development. In
particular, the analytic Stein solution used in \cref{supp:pivot:lem:palm} remains
an input to its Lean proof. The finite coupling argument is formally checked
under this premise, while the semigroup construction given in the companion is
outside that formal proof. This is a limitation of formal coverage, not an
additional assumption in the written argument.

The formal proofs contain no proof placeholders and use only Lean's standard
foundational axioms \texttt{propext}, \texttt{Classical.choice} and
\texttt{Quot.sound}. This foundational check does not discharge the literature
hypotheses above. The written proofs and the formal development provide
complementary accounts, whose precise correspondence is documented in the
repository.

\Needspace{17\baselineskip}
Seven Palomar registrations cover a selection of 89 declarations from the
development. All entries below are version~1; the counts refer to the registered
selection, not to the total number of results in the repository.

\begin{center}
\small
\setlength{\tabcolsep}{4pt}
\renewcommand{\arraystretch}{1.12}
\begin{tabularx}{\textwidth}{@{}>{\raggedright\arraybackslash}X r l@{}}
\toprule
\textbf{Scope} & \textbf{Count} & \textbf{Palomar identifier} \\
\midrule
Critical Poisson fields and conditional transfers & 16 &
\href{https://palomar-registry.org/entry?id=PALOMAR-2026-09-20-000003\&version=1}{\texttt{PALOMAR-2026-09-20-000003}} \\
Word dictionaries, exact marks and compound clusters & 11 &
\href{https://palomar-registry.org/entry?id=PALOMAR-2026-09-20-000004\&version=1}{\texttt{PALOMAR-2026-09-20-000004}} \\
Threshold staircases and the Poisson--Gaussian bridge & 29 &
\href{https://palomar-registry.org/entry?id=PALOMAR-2026-09-20-000007\&version=1}{\texttt{PALOMAR-2026-09-20-000007}} \\
Microscopic boundary, prefixes and longest runs & 7 &
\href{https://palomar-registry.org/entry?id=PALOMAR-2026-09-20-000011\&version=1}{\texttt{PALOMAR-2026-09-20-000011}} \\
Microscopic--bulk crossover and conditional affine laws & 11 &
\href{https://palomar-registry.org/entry?id=PALOMAR-2026-09-20-000012\&version=1}{\texttt{PALOMAR-2026-09-20-000012}} \\
Microscopic signed fields and empirical Poisson laws & 6 &
\href{https://palomar-registry.org/entry?id=PALOMAR-2026-09-21-000002\&version=1}{\texttt{PALOMAR-2026-09-21-000002}} \\
Regular configurations, Palm deficits and cumulant obstructions & 9 &
\href{https://palomar-registry.org/entry?id=PALOMAR-2026-09-21-000003\&version=1}{\texttt{PALOMAR-2026-09-21-000003}} \\
\bottomrule
\end{tabularx}
\end{center}

Each registration identifies a precise set of formal statements and a fixed
repository revision. Comparator checks agreement between the selected statements
and their proved counterparts; the proofs are checked by Lean and independently
replayed through NanoDa. Palomar's separate editorial review assesses the
correspondence with the informal account and the mathematical interest of the
submission. That review is automated and does not constitute independent human
peer review. The registrations document the verification of the selected formal
results under their stated hypotheses; they do not certify every argument or
claim in the article and companion. The verification and review protocol is
specified in
\href{https://github.com/PalomarRegistry/PalomarPolicy/blob/main/docs/specification.md}{\texttt{PalomarPolicy}}.
